% !TEX root = ../main.tex

\documentclass[../main.tex]{subfiles}

\graphicspath{{\subfix{../images/}}}

\begin{document}

\subsection{Внешний вид программы}
 
На рисунке~\ref{fig:main} представлен общий вид окна визуализатора в
реальном времени. Верхняя область содержит 3D-сцену с
моделью полётного контроллера на фоне координатной сетки и мировых осей;
нижняя --- четыре графика датчиков.
 
\begin{figure}[H]
    \centering

    \includegraphics[width=\linewidth]{images/screen_main.jpg}
    
    \caption{Главное окно визуализатора: 3D-сцена и HUD-панель}
    \label{fig:main}
\end{figure}
 
На рисунке~\ref{fig:rotate} показано состояние визуализатора при вращении
полётного контроллера: 3D-модель следует за реальной ориентацией устройства,
а кривые гироскопа на HUD-панели фиксируют угловые скорости по каждой из осей.
 
\begin{figure}[H]
    \centering
   
    \includegraphics[width=\linewidth]{images/screen_rotate.jpg}
    
    \caption{Визуализатор при вращении устройства: реакция HUD-графиков гироскопа}
    \label{fig:rotate}
\end{figure}
 
На рисунке~\ref{fig:hud} приведён укрупнённый вид HUD-панели. Нулевые линии
в каждом графике выделены серым. Видно, что в статике (покой платформы) кривые
гироскопа сходятся к нулю~--- это подтверждает корректность подобранных
смещений~\eqref{eq:biases}.
 
\begin{figure}[H]
    \centering
    
    \includegraphics[width=\linewidth]{images/screen_hud.jpg}
    
    \caption{Укрупнённый вид HUD-панели: акселерометр, гироскоп, углы Эйлера, скорость}
    \label{fig:hud}
\end{figure}
 

\subsection{Результаты оценки ориентации}
 
В ходе тестирования подтверждена корректная работа алгоритма кватернионного
интегрирования. При вращении полётного контроллера по всем трём осям
3D-модель точно воспроизводила движения устройства. Углы Эйлера, вычисленные
по формулам~\eqref{eq:roll}--\eqref{eq:yaw}, соответствовали физическим
поворотам.
 
\subsection{Результаты оценки позиции}
 
Алгоритм двойного интегрирования ускорения показал неудовлетворительные
результаты при оценке абсолютных координат. Ключевая проблема~---
\textbf{дрейф позиции}: уже через 5--10~с накопленная ошибка достигала
нескольких метров. Механизм ZUPT снизил скорость нарастания дрейфа в
состоянии покоя, однако не устранил его полностью из-за ненулевого шума
акселерометра порядка $10^{-3}$--$10^{-2}$~м/с$^2$.
 
\subsection{Анализ причин дрейфа}
 
Пусть истинное ускорение $a(t)$ и измеренное $\hat{a}(t)$ связаны:
\begin{equation}
    \hat{a}(t) = a(t) + b + \eta(t),
    \label{eq:noise}
\end{equation}
где $b$~--- остаточное смещение после калибровки, $\eta(t)$~--- случайный
шум. Ошибка позиции после двойного интегрирования за время $T$:
\begin{equation}
    \delta p(T) = \tfrac{1}{2}\,b\,T^2 +
    \int_0^T\!(T-t)\,\eta(t)\,dt.
    \label{eq:drift}
\end{equation}
 
Первое слагаемое растёт \textit{квадратично} по времени даже при малом
остаточном смещении~\cite{groves2013}. За 10~с это даёт ошибку порядка
нескольких метров~\cite{woodman2007}.
 
\subsection{Возможные направления развития}
 
Полноценная инерциальная навигация без внешних корректирующих источников
принципиально невозможна~\cite{groves2013}. Среди возможных улучшений:
\begin{enumerate}[leftmargin=2cm]
    \item \textbf{Фильтр Калмана с GPS-коррекцией}~\cite{groves2013}~---
          классический подход для авиационных INS.
    \item \textbf{Фильтр Маджвика}~\cite{madgwick2011}~--- градиентный
          алгоритм слияния данных гироскопа, акселерометра и магнитометра.
    \item \textbf{Нелинейный дополняющий фильтр на SO(3)}~\cite{mahony2008}.
    \item \textbf{ZUPT с фильтром Калмана}~\cite{skog2010}~---
          расширенное применение метода нулевой скорости.
\end{enumerate}
 
\newpage
	
\end{document}
